\chapter{Roadmap, Open Gaps, and How to Re-Audit This Book}
\index{roadmap}\index{re-audit}\index{open gap}
\label{ch:roadmap}

This roadmap is a snapshot against CNA \texttt{7a64362e} and the sibling revisions pinned in the
edition plan.  It deliberately avoids one global completion percentage.  API presence, behavior,
renderer evidence, platform execution, and automation are different axes, and several CNA
documents carrying headline percentages are already stale about XNB, glTF, renderers, or tests.

\section{Triage correctness before expanding breadth}

The source audit recorded 54 CNA findings in \texttt{cnabugs.md}.  They were derived by static
reading and configuration analysis; CNA itself was not built or modified during that audit.  The
first upstream step is therefore triage and focused reproduction, not bulk patching.

High-value starting points have small counterexamples and broad consequences: aliased
\texttt{Matrix::Transpose} corrupts \texttt{Plane::Transform}; StorageContainer lacks path
containment; the CNJ parser has no depth bound; the main CI workflow cannot configure; several
capability queries promise operations their renderer refuses.  Each should receive a minimal
failing test, a scoped fix, controls on unaffected renderers, and a current evidence record.

The audit list must not become a second stale bug tracker.  Once triaged upstream, rows need an
issue or commit pointer and a status tied to a reproducer.  Findings rejected as intentional
behavior should move into the documented-deviation contract rather than disappear.

\section{Complete capability truth before adding renderers}

CNA already exposes \RendererIdentityCount{} identities across \RendererFamilyCount{} families.
The next renderer milestone is not a
larger count; it is truthful answers from the existing contract.  HEADLESS, SOFTWARE, and WEBGPU
advertise MRT while rejecting multi-target use; WebGPU and bgfx overclaim custom effects; several
query and stencil reports disagree with implementation; DIRECTX10 inherits an unconditional
colored extended-draw downgrade.

Build a table-driven capability conformance suite that pairs every \texttt{true} with a usable
factory/draw case and every \texttt{false} with a deterministic refusal.  Keep capability tests
renderer-owned where one configure can exercise them and retain cross-renderer public fixtures as
controls.  A registry count cannot substitute for this behavioral matrix.

\section{Close the arbitrary-effect boundary deliberately}

The public compiled-effect constructor and general XNB \texttt{EffectReader} remain unsupported.
FNA3D proves that real XNA stock \texttt{.fxb} blobs can execute through MojoShader, while D3D9
vendors source and compares much of its compiled instruction stream.  These are strong foundations,
not a completed arbitrary-effect pipeline.

A roadmap decision must define accepted bytecode, reflection, parameter and technique mapping,
renderer translation routes, failure behavior, and licensing/provenance.  Vulkan needs a path from
MojoShader's output to SPIR-V; bgfx expects offline binaries; several renderers intentionally have no
custom-effect support.  Do not hide that matrix behind one Boolean or widen the public constructor
until each supported identity has an engagement test.

\section{Harden content symmetrically}

XNB support is real and substantial.  Remaining work includes whole-file size enforcement, named
unsupported reader/version policy, LZ4's explicit deferred status, writer/tooling decisions, and
continued hostile-input tests.  CNJ needs the same depth/size/fuzz discipline already applied to
XNB; its current recursive parser is an obvious first target.

The façade has additional contract work: \texttt{Unload()} disposes nothing,
\texttt{CanDeserializeIntoExistingObject} is never consulted, Song/Video can resolve a same-name
CNJ they cannot parse, and manifest data is introspection rather than resolution.  Fixes should
retain the separation between caller-selected asset paths and file-internal reference containment.

\section{Complete glTF semantics beyond importer coverage}

The conformance ladder proves D1--D4 and D8 repairs and exposes remaining material and animation
limits.  Runtime import still discards a computed PBR UV-set mismatch warning, ignores declared
required extensions, supports only the first mesh group in one route, and applies position-only
morphing to some PBR strides without diagnosis.  D6/D7-era material and rigid-animation work
remain the clearest compatibility targets.

Advance one independently generated fixture at a time.  Update the manifest, scene digest,
semantic golden, runtime/offline parity, and renderer evidence together.  Run a pinned external
glTF validator rather than claiming generator self-validation that the current tool does not
perform.

\section{Turn platform claims into retained artifacts}

The adapted Android pin needs a fresh graphics build and emulator frame before it receives a
current renderer claim.  Retain
the APK, configure log, logcat, screenshot, ABI/API/NDK identity, and lifecycle trace.  Then cover
background, foreground, low-memory, termination, asset access through ContentManager, and physical
hardware.

The adapted Metal tree needs a native compile, link, validation, pixel, Retina, and frame-pacing run;
the historical pre-adaptation workflow does not close that boundary.  iOS/tvOS require a designed
host and remain explicit non-claims.  WebGL1, WebGL2, and Canvas need the HTTP-served browser verdict
and compositor pattern already demonstrated by HTML-DOM/SVG-DOM, plus context-loss and shader
compilation evidence.

For Windows, retain the MinGW/Wine route while adding genuine native MSVC evidence where the public
support table claims it.  Every result must name translation runtime, GPU/software device, display,
and prefix.  A compatibility-layer pass is valuable and is not a native-Windows pass.

\section{Repair the verification apparatus itself}

Replace obsolete \texttt{EASYGL} workflow selectors with public identities and make the full Linux
and sanitizer legs reach compilation.  Update the four post-modularization validator paths, run the
Direct2D mutation campaign rather than only anchor dry-run, fix Device filter drift, and decide one
authoritative direct-binary-versus-CTest policy for the full suite.

Label all renderer registrations consistently or publish generated selectors that do not depend on
labels.  Make link-closure evidence generator-independent instead of silently skipping under the
project's own presets.  Bring SVG-DOM's existing browser runner into automation and ensure the
release gates, not only their non-release validators, are exercised.

Add code-coverage tooling only with a defined question.  Line coverage can identify untouched code
but cannot replace oracle authority, mutation sensitivity, or runtime engagement.  Keep API-surface
``coverage'' under a different name so the metrics cannot be confused.

\section{Make CNA consumable as a package}

The pinned build is designed around sibling \texttt{add\_subdirectory} composition and has no
complete install/export package for CNA.  A mature consumer path should define installed module
targets, public headers, renderer selection, dependency discovery, version compatibility, and a
small out-of-tree package consumer test.  Sibling libraries such as meta-gl already demonstrate
that installed-package and exported-symbol tests can make this boundary executable.

Packaging work must preserve the one-SDL-instance invariant and optional-module closure.  It should
not force every renderer dependency onto a consumer that selects one identity.

\section{Reconcile documentation from generated facts}

Remove or banner claims that CNA cannot read XNB, that WebGL2 or Android graphics is pixel-verified,
that the implementation-family count means public identities, or that old source paths still register gates.  Generate renderer
identity/family tables, reader inventories, workflow summaries, and test-selection tables from the
same sources their validators use.

Dated reports may remain for history if their revision and superseding authority are prominent.
Comments should explain why; scripts and tests should establish whether.  The fastest way to create
new drift is to copy a corrected number into five prose files without a derivation.

\section{How to re-audit this edition}

A future edition should repeat a bounded protocol:

\begin{enumerate}
  \item pin CNA and every cited sibling SHA; record branch divergence and keep them read-only;
  \item enumerate physical modules, public renderer identities/families, build options, workflows,
        scripts, test registration paths, content readers, and sample/example catalogs mechanically;
  \item re-run focused source audits by subsystem and record contradictions separately from defects;
  \item verify every explicit path and identifier in the manuscript against the pin;
  \item rebuild the complete book, require clean references/index/stale-term checks, render every
        changed physical page, and inspect high-risk tables and listings at full size;
  \item update \texttt{PLAN.md} and \texttt{NEXT.md} with commands, counts, page ranges, visual
        findings, fixes, and the next bounded batch.
\end{enumerate}

The edition should change when evidence changes, not merely when prose ages.  The final lesson of
CNA's own verification culture applies to this book as well: a claim is durable only when a future
reader can prove that the intended source, runtime, oracle, and page actually engaged.
